%! Author = chtompki
%! Date = 1/14/26
% Example LaTeX document for GP111 - note % sign indicates a comment
\documentclass[12pt]{amsart}
% Default margins are too wide all the way around. I reset them here
\setlength{\hoffset}{-.75in}
\setlength{\textwidth}{6.5in}
\setlength{\voffset}{-.5in}
\setlength{\textheight}{9.0in}
\usepackage{amsmath}
\usepackage{amssymb}
\usepackage{amsthm}
\usepackage{pgfplots}
\usepackage{enumerate}
\usepackage{tikz}
\usetikzlibrary{shadows,arrows,arrows.meta,shapes,positioning,calc}% Required for the Circle and Triangle arrow tips
\usepackage{setspace}
\usepackage{siunitx}
\usepackage{enumitem}

\theoremstyle{conjecture}
\newtheorem{conjecture}{Conjecture}

\theoremstyle{definition}
\newtheorem{definition}{Definition}

\theoremstyle{question}
\newtheorem{question}{Question}

\theoremstyle{theorem}
\newtheorem{theorem}{Theorem}

\newenvironment{solution}
{\begin{proof}[Solution]}
{\end{proof}}

\title{Graph Theory:\\Homework 5\\ Book Section 3.1.}
\author{Rob Tompkins}
\date{\today}


\begin{document}

    \maketitle
    \noindent\textbf{1.} let $S$ be the set of vertices saturated by a matching $M$ in a graph $G$.
    Prove that some maximum matching also saturates all of $S$.
    (Hint: Use Berge's Theorem to show that on can start with any matching $M$, and find a sequence of
    ``alternating path augmentations'' that produce a maximum matching $M^*$.
    Be careful not to conflate ``maximum'' with ``maximal.'')

    \begin{proof}
    \end{proof}

    \noindent\textbf{2.} Determine the minimum size of a maximal matching in the cycle $C_n$,

    \begin{solution}
    \end{solution}

    \noindent\textbf{3.} Let $M$ and $M'$ be matchings in an $X,Y$-bigraph $G$.
    Suppsose that $M$ saturates $S\subseteq X$ and $M'$ saturates $T\subseteq Y$.
    Prove that $G$ has a matching that saturates $S \cup T$

    \begin{proof}
    \end{proof}

    \noindent\textbf{4.} When $k\ge 2$, prove that $Q_k$ has at least $2^{2^{k-2}}$ perfect matchings

    \begin{proof}
    \end{proof}

    \noindent\textbf{5.} Let $G$ be the Petersen Graph.
    \begin{itemize}
        \item[(a)] Prove that deleting any perfect matching of $G$ leaves $C_5 + C_5$ (i.e. two disjoint 5-cycles).
        (Hint: After you remove a perfect matching what is the degree of each vertex).
        \item[(b)] Prove that the Petersen graph has exactly six perfect matchings.
        (Hint: use (a) and use the fact that you already know how many 5-cycles $G$ contains.)
    \end{itemize}

    \begin{proof}
    \end{proof}


\end{document}